\documentclass{ximera}

\title{u-Substitution Review Activity}
\author{MATH 426: Calculus II}

\begin{document}
\begin{abstract}
   Working with the peers in your group, solve the following problems. Make sure to show and justify all your work. Make sure everyone in the group understands the solution and participates. Be prepared to report your answers to the whole class. 
\end{abstract}
\maketitle

The start of this activity is more open-ended and conceptual than the other activities we have seen so far in the course. As we move forward with the content, you will be exposed to several independent techniques of integration. To fully master this material you will need to be able to \textbf{choose} appropriate techniques for problems `in the wild.' This choice is \textbf{important} and it is not something that we show in our written work. In our first instantiation of this activity you will be asked to discuss your strategy for selecting an integration technique, and to share your insights with your peers. If you can begin to articulate what you `see' in mathematical problems you will be better prepared to choose the correct procedures for handling new problems.\\


\begin{exercise}
    Discuss the different integration techniques we have discussed so far in Calculus 2 (FTC and $u$-substitution). (Consider the following questions as discussion starters, you need not address all the questions in your discussion, but try to hash out any insight or confusion using these questions as a guide.)

    \begin{itemize}
        \item What makes an integration problem easy (or hard)?

        \item What features of the integrand make you think of specific techniques?

        \item If you are unsure how to proceed on a problem, what do you do next?

        \item Do you treat definite and indefinite integrals differently? How?

        \item What kinds of problems can you do in your head?
        
        \item After your discussion, compress your thoughts into a single piece of advice which is helpful for you:
    \end{itemize}
\end{exercise}

\begin{exercise}
    As a group create one example problem that you think is the perfect example of each of the following situations: (You can create group examples, or choose your own personal examples. You do not need to complete the problems here, but it may be helpful to annotate the problems by adding notes showing the features you see.)
        \begin{enumerate}
            \item An integral which relies on the fundamental theorem of calculus.
            
            \item An integral which has a \textbf{perfect} $u$-substitution.
            
            \item An integral which has an \textbf{imperfect} $u$-substitution.
            
            \item An integral which requires an \textbf{affine (or linear)} substitution to be tractable.
        \end{enumerate}
\end{exercise}

\begin{exercise}
    Evaluate the following integrals.
    \begin{enumerate}
        \item $\int \cot(x)dx $ (hint: you will need to rewrite $\cot(x)$ in order to be able to complete this integral.)
        \item $\int_0^1 \frac{e^x}{5+e^x}dx$
        \item $\int x^2(x^3+4)^6dx$
        \item $\int \frac{(\ln(x)^5)}{x}dx$
        \item $\int_{-1}^2 2x^3(x^4+3)dx$
    \end{enumerate}
\end{exercise}

\begin{exercise}
    Evalute the integrals (a) $\int x\sqrt{x-1}dx$, (b) $\int x\sqrt{x^2-1}dx$ and (c) $\int x^2\sqrt{x-1}dx$.  How are the methods for each of these integrals the same?  How are they different?  What characteristics caused you to choose your method for each integral?
\end{exercise}



\end{document}